% !TeX encoding = UTF-8
% !TeX spellcheck = pl_PL

\def\filename{wstep}
\chapter{Wstęp}
\section{Motywacja}
Dynamiczny rozwój technologii sztucznej inteligencji w ostatnich latach skłania środowisko naukowe do intensywnego badania jej zastosowań w różnych obszarach medycyny, w szczególności w diagnostyce obrazowej.
Konwolucyjne sieci neuronowe wykazały znaczący potencjał w zakresie komputerowo wspomaganej diagnostyki, osiągając wysoką skuteczność w analizie obrazów medycznych.
Motywacją niniejszej pracy jest pogłębienie wiedzy na temat możliwości wykorzystania konwolucyjnych sieci neuronowych w diagnostyce zmian skórnych oraz opracowanie prototypu mobilnego systemu wspierającego wstępną ocenę znamion.
System ten ma na celu zarówno zwiększenie dostępności narzędzi diagnostycznych, jak i podniesienie świadomości pacjentów w zakresie monitorowania zdrowia skóry.

\section{Cele pracy}
Celem pracy jest opracowanie prototypu systemu rozpoznawania zmian skórnych opartego na konwolucyjnych sieciach neuronowych, zdolnego do identyfikacji rodzaju zmiany skórnej uchwyconej na zdjęciu wykonanym aparatem telefonu przez użytkownika.
Projektowana aplikacja powinna charakteryzować się wysoką wydajnością oraz umożliwiać inferencję w czasie rzeczywistym, tak aby użytkownik mógł szybko przeprowadzić wstępną analizę niepokojących zmian skórnych.

\section{Metodyka badawcza}
W ramach badań przedstawiono system klasyfikacji zmian skórnych wykorzystujący konwolucyjną sieć neuronową zaimplementowaną w postaci aplikacji mobilnej. 
Wykorzystano zbiór PAD-UFES-20, obejmujący zdjęcia wykonane aparatem telefonu komórkowego, na podstawie których dotrenowano siedem wybranych architektur.
Najlepszy model został następnie przekonwertowany do formatu mobilnego i poddany wstępnym badaniom empirycznym na rzeczywistych zdjęciach zmian skórnych.

\section{Struktura pracy}
Praca składa się z pięciu rozdziałów.
Rozdział drugi przedstawia potencjał sztucznej inteligencji oraz urządzeń mobilnych w obszarze diagnostyki wspomaganej komputerowo, omawia dotychczasowe rozwiązania w tej dziedzinie oraz podstawy głebokiego uczenia i metody oceny skuteczności modeli.
Uwzględniono również kluczowe wyzwania związane z wykorzystaniem AI w diagnostyce dermatologicznej.
Rozdział trzeci opisuje dostępne zbiory danych dotyczące zmian skórnych oraz proces decyzyjny prowadzący do wyboru modelu wykorzystanego w systemie.
Przedstawiono również charakterystykę zmian skórnych uwzględnionych w badaniu oraz teoretyczną infrastrukturę projektowanego systemu.

W rozdziale czwartym zaprezentowano szczegółową metodykę badawczą, obejmującą przygotowanie danych, proces trenowania modeli oraz implementację aplikacji mobilnej na system Android, wraz z diagramem klas ilustrującym architekturę aplikacji i jej komunikację z modelem.
Przeanalizowano wyniki uzyskane dla wszystkich trenowanych modeli, ich skuteczność w analizie poszczególnych kategorii zmian skórnych oraz rezultaty badań empirycznych przeprowadzonych na łagodnych zmianach skórnych.
